\section{Relative primary signatures and specialization laws}
\label{sec:relative-primary-specialization}

This section supplies the two family arguments that are not consequences
of a single exact witness: genericity on coefficient orbits and
fibrewise control of associated points.  It also extracts the
geometric primary signatures of the corank-two layers and writes
explicit orbit degenerations.

\subsection{Stabilizer slices and generic Groebner bases}

Keep
\[
 P=\C[a,b,c,d],\qquad
 T=\begin{pmatrix}a&b\\c&d\end{pmatrix},\qquad
 \delta=ad-bc.
\]
For a mixed normal form \(\beta\), let
\(G_\beta\subset GL(H)\times GL(W)\times GL(E_H)\) be its
stabilizer.  If \((X,Y,Z)\) lies in its Lie algebra, then
\begin{equation}\label{eq:stabilizer-linearization}
 Z\beta-\beta(X\otimes1+1\otimes Y)=0,
\end{equation}
and its infinitesimal action on the quadratic block is
\begin{equation}\label{eq:chi-linearization}
 \dot\chi=Z\chi-\chi\,\Sym^2Y.
\end{equation}
These are linear equations over \(\Q\) at the integral base
points of Lemma~\ref{lem:nine-orbit-colons}.

\begin{theorem}[Generic slice certification for the mixed-kernel atlas]
\label{thm:generic-slice-certification}
For each of the seven mixed-kernel rows not already covered by the
mixed-regular and decomposable-kernel theorems, the Hilbert functions
in Theorem~\ref{thm:corank-two-mixed-kernel-atlas} are the generic
Hilbert functions on the entire orbit family.

More precisely, at the integral base points used in
Lemma~\ref{lem:nine-orbit-colons}, the ranks of the infinitesimal
\(G_\beta\)-orbits in the nine-dimensional \(\chi\)-space and
one choice of transverse affine slice are
\[
\begin{array}{c|c|c}
\text{type}&\rank T_{\chi}(G_\beta\chi)&\text{slice entries}\\ \hline
\text{secant}&6&\chi_{11},\chi_{21},\chi_{31}\\
\text{tangent}&7&\chi_{11},\chi_{21}\\
H\text{-ruling}&7&\chi_{11},\chi_{13}\\
W\text{-ruling}&9&\varnothing\\
\text{dual rank }2&9&\varnothing\\
\text{dual rank }1&9&\varnothing\\
\text{zero mixed map}&9&\varnothing.
\end{array}
\]
On the first three slices, with coordinates \(s_i\), Groebner
bases of all required colons over \(\Q(s_i)\) give respectively
\[
\begin{array}{c|c|c|c}
\text{type}&\min\{k:\delta^k\in J\}&H_{W_2}&H_{W_3}\\ \hline
\text{secant}&2&(1,4,3)&0\\
\text{tangent}&2&(1,4,3)&0\\
H\text{-ruling}&3&(1,4,6,4)&(1).
\end{array}
\]
The only nonconstant denominators in reduced bases may be chosen to
divide
\[
 s_2-2\quad\text{(secant)},\qquad
 s_0+2s_1-5\quad\text{(\(H\)-ruling)},
\]
and neither vanishes at the base point \(s=0\); the tangent slice
requires no nonconstant denominator.  In the remaining four rows the
stabilizer orbit is already open in \(\chi\)-space, so the exact
base-point calculation is the generic calculation.
\end{theorem}

\begin{proof}
Solving \eqref{eq:stabilizer-linearization} and applying
\eqref{eq:chi-linearization} gives the ranks in the table.  Adjoining
the displayed coordinate directions raises the tangent rank to nine.
Thus the morphism from the stabilizer times the affine slice to
\(\chi\)-space has surjective differential at the base point.
Over characteristic zero it is therefore dominant and smooth on a
neighbourhood of that point.

Compute the residual minors and successive colons in
\(\Q(s_i)[a,b,c,d]\).  For the secant and tangent slices the
second-layer generic initial ideal is
\[
 (c,d)^3+(a,b)^2+(a,b)(c,d),
\]
which has Hilbert function \((1,4,3)\), and the next colon is the
unit ideal.  For the \(H\)-ruling slice a generic initial ideal of
the second layer is generated by
\[
 c^4,c^3d,c^2d^2,cd^3,d^4,
 ac^2,acd,ad^2,bd^2,a^2,ab,b^2,bc,
\]
with Hilbert function \((1,4,6,4)\); the third-layer initial ideal
is \((a,b,c,d)\).  These initial ideals are unchanged after
specialization away from the displayed denominator divisors.

For the four open-orbit rows the exact Groebner bases give the
functions printed in Theorem~\ref{thm:corank-two-mixed-kernel-atlas}.
The calculation is reproduced, over exact rational function fields
rather than floating-point samples, by
\texttt{checks/generic\_boundary\_atlas.py}.  Since the slice map
is dominant, these are the generic initial ideals on the orbit
families.
\end{proof}

\subsection{Geometric primary signatures for the nine mixed-kernel types}

Let \(Q=V(\delta)\) be the affine Segre cone and let
\(\mathfrak m=(a,b,c,d)\) be its vertex.  Write \(R_r\) for a
reduced divisor consisting of \(r\) distinct lines in one ruling of
\(Q\).  Write \(V_\ell\) for a homogeneous
\(\mathfrak m\)-primary scheme of length \(\ell\).  Finally,
\(R_r+E_\ell\) means a scheme whose saturation is \(R_r\) and
whose saturation quotient is an embedded vertex module of length
\(\ell\).  This notation records associated supports and generic
lengths; it does not choose a noncanonical embedded primary
component.

\begin{theorem}[Generic primary signatures for the mixed-kernel orbits]
\label{thm:generic-primary-signatures}
On the generic coefficient stratum of each \(A_H\)-injective
mixed-kernel orbit, the positive layers have the following primary
signatures:
\[
\begin{array}{c|c|c}
\text{mixed-kernel type}&W_2&W_3\\ \hline
3,\ \text{point off }Q&V_1&\varnothing\\
3,\ \text{point on }Q&V_3&\varnothing\\
2,\ \text{secant}&V_8&\varnothing\\
2,\ \text{tangent}&V_8&\varnothing\\
2,\ H\text{-ruling}&V_{15}&V_1\\
2,\ W\text{-ruling}&R_1+E_5&\varnothing\\
1,\ \text{dual rank }2&R_2+E_3&V_1\\
1,\ \text{dual rank }1&R_3&V_3\\
0&Q&Q.
\end{array}
\]
In particular, every associated support and every generic length in
these layers is determined.  The symbols \(R_r\) have \(r\)
minimal ruling-line primes, each of generic length one.  The only
additional associated prime in \(R_1+E_5\) and
\(R_2+E_3\) is the embedded vertex \(\mathfrak m\).
\end{theorem}

\begin{proof}
The finite rows are homogeneous zero-dimensional schemes supported at
the vertex, hence are \(\mathfrak m\)-primary; their lengths are
the Hilbert-function sums.  For the \(W\)-ruling witness the second
layer is
\[
 I=(ac,bc,ad,bd,c^3,c^2d,cd^2,d^3).
\]
Its saturation is \((c,d)\), and \((c,d)/I\) has length five.
Thus its associated primes are \((c,d)\) and \(\mathfrak m\),
with an embedded vertex contribution of length five.

For the dual-rank-two witness, saturation gives
\[
 (a,b)\cap(c,d),
\]
the union of two ruling lines, and the saturation quotient has length
three.  For the dual-rank-one witness, modulo \(\delta\), the
second-layer ideal is exactly
\[
 \bigcap_{\lambda\in\{0,1,3\}}
 (a+\lambda c,b+\lambda d),
\]
so it is a reduced union of three ruling lines and has no embedded
point.  The third layer is the length-three vertex thickening already
computed in Lemma~\ref{lem:nine-orbit-colons}.  The zero mixed-map
identity \(J=(\delta^3)\) gives \(W_2=W_3=Q\).

Theorem~\ref{thm:generic-slice-certification} transports these exact
signatures to the generic coefficient strata.  Associated supports
are preserved by the stabilizer action, and the relative-primary
argument below prevents an isolated base-point calculation from being
used beyond its certified open stratum.
\end{proof}

\subsection{Pure-block rank drop at projection corank two}

When \(A_H\) has rank \(a<3\), put \(q=6-a\) and
\(b=\rank\overline B_H\).  For \(b=3,2,1\) the phrase
\emph{generic mixed kernel} means respectively an off-Segre kernel
point, a secant kernel line, and a kernel plane whose dual tensor has
rank two.

\begin{theorem}[Generic primary signatures for every \(A_H\)-rank-drop type]
\label{thm:rankdrop-primary-signatures}
At projection corank two, every generic rank type allowed by
Proposition~\ref{prop:rankdrop-coranktwo-depth} has the following
complete list of nonzero positive layers:
\[
\begin{array}{c c|c|l}
a&b&\text{nilpotency index}&(W_1,W_2,W_3,W_4,W_5)\\ \hline
2&4&3&(Q,V_{22})\\
2&3&4&(Q,V_{27},V_1)\\
2&2&4&(Q,R_4,V_8)\\
2&1&5&(Q,Q,Q,V_1)\\
1&4&5&(Q,Q,V_{11},V_1)\\
1&3&5&(Q,Q,R_2+E_3,V_1)\\
1&2&5&(Q,Q,Q,R_2)\\
0&4&6&(Q,Q,Q,R_2,V_1)\\
0&3&6&(Q,Q,Q,Q,V_1).
\end{array}
\]
Here every \(Q\) is the reduced Segre quadric with generic length
one; every \(R_r\) is a reduced union of \(r\) ruling lines with
generic length one on each component; and every \(V_\ell\) is
\(\mathfrak m\)-primary of the indicated length.  In the unique
row \((a,b)=(1,3)\) carrying \(E_3\), the associated primes of
that layer are the two ruling-line primes together with the embedded
vertex.
\end{theorem}

\begin{proof}
Choose the generic normal forms of the mixed block described above and
let \(\chi\) vary.  The stabilizer-slice codimensions in
\(\chi\)-space for
\[
 (a,b)=(2,4),(2,3),(2,2),(2,1),(1,4),(1,3),(1,2),(0,4),(0,3)
\]
are respectively
\[
 5,4,1,0,3,1,0,0,0.
\]
On explicit complementary coordinate slices the successive colons are
computed over the rational function fields of those slices.  Their
Hilbert functions are
\[
\begin{array}{c c|l}
a&b&\text{non-quadric Hilbert functions}\\ \hline
2&4&(1,4,9,8)\\
2&3&(1,4,9,8,5);\ (1)\\
2&2&4n+4\ (n\ge4);\ (1,4,3)\\
2&1&(1)\\
1&4&(1,4,6);\ (1)\\
1&3&2n+2\ (n\ge3);\ (1)\\
1&2&2n+2\ (n\ge1)\\
0&4&2n+2\ (n\ge1);\ (1)\\
0&3&(1).
\end{array}
\]
Every omitted earlier layer has Hilbert function \((n+1)^2\),
hence equals \(Q\).

To identify the positive-dimensional supports, pull the layer ideals
back to the rank-one parametrization \(T=uv^t\).  At integral base
points the common binary factors are
\[
\begin{array}{c c|c}
a&b&\text{binary factor in }u\\ \hline
2&2&u_0u_1(9u_0^2+6u_0u_1-7u_1^2)\\
1&3&10u_0^2+u_0u_1-4u_1^2\\
1&2&u_0u_1\\
0&4&6u_0^2+u_0u_1+4u_1^2.
\end{array}
\]
All have distinct roots.  The first therefore cuts four reduced
ruling fibres and the remaining three cut two.  Their diagonal Hilbert
functions on the Segre cone are exactly \(4n+4\) from degree four
and \(2n+2\) from degree two, respectively.  In the \((1,3)\)
row the actual Hilbert function exceeds the reduced two-line divisor
by three in degree two and agrees thereafter; saturation therefore
leaves a length-three vertex module, giving \(R_2+E_3\).

The exact stabilizer ranks, slice coordinates, Groebner leading
monomials, denominator checks, and Hilbert counts are reproduced by
\texttt{checks/generic\_boundary\_atlas.py}.  Because each slice
map is dominant at its base point, the calculation is at the generic
point of the rank type rather than at a special witness.
\end{proof}

\subsection{Relative assassins and base change for the colon tower}

We now prove Proposition~\ref{prop:associated-prime-refinement}.
The point is to establish base change for the colons before making any
statement about fibrewise associated points.

\begin{proof}[Proof of Proposition~\ref{prop:associated-prime-refinement}]
Let \(S\) be a finite-type coefficient chart and put
\(X=\Spec P_S\), where \(P_S\) is the polynomial algebra in the
entries of \(T\).  For \(j\le5\), multiplication by
\(d^{j-1}\) gives an exact sequence of coherent
\(P_S\)-modules
\begin{equation}\label{eq:colon-kernel-family}
 0\longrightarrow K_j\longrightarrow P_S
 \xrightarrow{\ d^{j-1}\ } P_S/J
 \longrightarrow C_j\longrightarrow0,
\end{equation}
where \(K_j=(J:d^{j-1})\) and
\(C_j=P_S/(J,d^{j-1})\).

By generic flatness, followed by Noetherian induction on the complement,
there is a finite locally closed stratification of \(S\) on which
every \(C_j\) is flat over the base.  Tensoring
\eqref{eq:colon-kernel-family} with a residue field is then left
exact at \(P_S\):
\(\Tor_1^S(C_j,k(s))=0\).  Hence
\[
 (K_j)_s=(J_s:d_s^{j-1})
\]
for every geometric point of the stratum.  Refine once more so that
each
\(\mathcal M_j=P_S/((d)+K_j)\) and every cokernel of a graded
multiplication map is flat.  This proves base change for the graded
layers and constancy of their Hilbert polynomials and multiplication
ranks.

Fix one resulting noetherian stratum.  For the coherent flat module
\(\mathcal M_j\), consider its relative assassin in the sense of
\cite[Tag 05AS]{StacksRelativeAssassin}.  Over the generic point of
each irreducible component of the base there are finitely many
associated points.  After shrinking, their closures
\(Z_{j,\nu}\subset X\) are flat over the base and have geometrically
constant relative dimension.  The flat relative-assassin results of
\cite[Tag 0GSF]{StacksFlatRelativeAssassin}, in particular
Tag 0GSJ, identify the generic points of the nonempty fibres of these
closures with fibrewise associated points on a dense open.  The
specialization criterion in \cite[Tag 05L2]{StacksPurity} controls
their closure under specialization.

The locus on which an additional fibrewise associated point occurs is
constructible.  Remove its closure from the dense open just obtained
and repeat on the noetherian complement.  Noetherian induction
terminates and gives a finite stratification on which the list of
fibrewise associated points is exactly the list of generic points of
the nonempty \((Z_{j,\nu})_s\).  Refining by the constructible
incidence relations among these finitely many supports makes
minimal-versus-embedded status and all inclusions constant.

Finally localize \(\mathcal M_j\) at the generic point of each
\(Z_{j,\nu}\).  Its finite length is the rank of a finite filtration
over the function field of that support.  Generic freeness makes the
successive quotients locally free after shrinking, hence the generic
length is constant; Noetherian induction handles the complement.
This proves all clauses of the proposition.

The construction depends on presentations and on choices of
stratification, but its fibrewise output is
\(\Ass(\mathcal M_{j,s})\), generic length, and multiplication in
the intrinsic graded algebra.  Two presentation atlases therefore
have a common refinement with identical intrinsic signatures.
\end{proof}

\subsection{Explicit orbit closures and graded specialization}

Let
\[
 e_1=h_1w_1,\quad e_2=h_1w_2,\quad
 e_3=h_2w_1,\quad e_4=h_2w_2
\]
be the standard basis of \(H\otimes W\).  The following families
are written before applying a general quadratic block \(\chi_t\).

\begin{theorem}[Corank-two orbit closure and specialization laws]
\label{thm:coranktwo-specialization-laws}
The closure relations needed to pass between all nine mixed-kernel
types are realized by the following one-parameter families:
\begin{enumerate}
\item rank three, off \(Q\) to on \(Q\):
\[
 \beta_t=
 \begin{pmatrix}0&1&0&0\\0&0&1&0\\-t&0&0&1\end{pmatrix},
 \qquad \ker\beta_t=\langle e_1+t e_4\rangle;
\]
\item secant to tangent:
\[
 \ker\beta_t=\langle e_1+e_4,\ e_2-t e_3\rangle,
\]
for which \(\delta|_{\Pj(\ker\beta_t)}=x^2+t y^2\);
\item secant to the two ruling families:
\[
 \langle e_1,e_2+t e_4\rangle\rightsquigarrow
 \langle e_1,e_2\rangle,\qquad
 \langle e_1,e_3+t e_4\rangle\rightsquigarrow
 \langle e_1,e_3\rangle;
\]
\item tangent to dual-rank-two:
\[
 \beta_t=
 \begin{pmatrix}0&1&-1&0\\0&0&0&t\\0&0&0&0\end{pmatrix};
\]
\item either ruling to dual-rank-one, by multiplying the second
nonzero row in the corresponding rank-two normal form by \(t\);
\item either rank-one type to the zero map, by replacing
\(\beta\) with \(t\beta\).
\end{enumerate}
For a general polynomial choice of \(\chi_t\), the generic fibre
and special fibre lie in the generic coefficient strata of the
indicated orbit types.  The corresponding primary signatures
specialize as follows:
\[
\begin{array}{l|l}
\text{degeneration}&\text{change of positive layers}\\ \hline
\text{off }Q\to\text{on }Q&V_1\to V_3\\
\text{secant}\to\text{tangent}&V_8\to V_8\\
\text{secant}\to H\text{-ruling}&V_8\to(V_{15},V_1)\\
\text{secant}\to W\text{-ruling}&V_8\to(R_1+E_5)\\
\text{tangent}\to\text{dual rank }2&V_8\to(R_2+E_3,V_1)\\
H\text{-ruling}\to\text{dual rank }1&(V_{15},V_1)\to(R_3,V_3)\\
W\text{-ruling}\to\text{dual rank }1&(R_1+E_5)\to(R_3,V_3)\\
\text{dual rank }2\to0&(R_2+E_3,V_1)\to(Q,Q)\\
\text{dual rank }1\to0&(R_3,V_3)\to(Q,Q).
\end{array}
\]
Thus ruling components and embedded vertex components are born in
actual one-parameter families of Fitting schemes.
\end{theorem}

\begin{proof}
The determinant of a tensor in each displayed kernel gives the stated
Segre incidence.  For example
\[
 \det\left(xI_2+y\begin{pmatrix}0&1\\-t&0\end{pmatrix}\right)
 =x^2+t y^2,
\]
which is secant for \(t\ne0\) and tangent for \(t=0\).  Likewise
\(\delta\) restricts to \(txy\) on each secant-to-ruling family.
The row-scaling families have the stated rank-one limits by inspection.

The residual maximal minors are polynomial in \(t\) and in the
coefficients of \(\chi_t\).  Each source and target generic stratum
is a nonempty open in its coefficient space.  Choosing
\(\chi_0\) in the target open and a general first-order perturbation
\(\chi_0+t\chi_1\) avoids the source bad divisors over
\(\C(t)\).  Hence the two fibres have the primary signatures in
Theorem~\ref{thm:generic-primary-signatures}.  Because the ideals
come from one polynomial family, the changes displayed in the table
are genuine specializations, not comparisons of unrelated normal
forms.
\end{proof}

There is an analogous pure-block degeneration.  In a fixed target
basis let
\[
 A_t=(f_1,f_2,t f_3):\Sym^2H\longrightarrow S_R
\]
and choose the remaining mixed and quadratic coefficients generally.
For \(t\ne0\), the induced mixed block is generically the rank-three
off-Segre type; at \(t=0\), \(A_H\) has rank two and the induced
mixed block has rank four.  Thus the first nontrivial layer specializes
from \(V_1\) to \(V_{22}\).  This gives an explicit
\(A_H\)-injective-to-rank-drop family, while the full table of
Theorem~\ref{thm:rankdrop-primary-signatures} describes the generic
pure-block boundary reached after further rank drops.
